\chapter{The Shape of a Spark Application}
\label{ch:shape}

\begin{chapterintro}
A query you would actually write --- a week of \texttt{trips} joined to a route
dimension, aggregated to an on-time rate --- runs in four minutes. This chapter
does not make it faster. It uses it as a specimen: you run it, open the Spark UI,
and learn to name every moving part --- the driver, the executors, the one job,
its stages, the shuffle between them, and the one task doing more than the rest.
By the end you can read a running job instead of watching a progress bar.
\end{chapterintro}

\section{A query, end to end}
\tbd

\section{Driver, executors, and the cluster manager}
\tbd

\section{Deploy modes and Spark on Kubernetes}
\tbd

\section{Jobs, stages, and tasks}
\tbd

\section{Narrow versus wide dependencies}
\tbd

\section{Reading the UI for the first time}
\tbd

\section{Summary}
\tbd

\exercisehead
\begin{exercises}
\item Run two queries against \texttt{trips}, one with a \texttt{GROUP BY} and one
      without. Identify which produced a second stage, and explain why.
\end{exercises}
